\documentclass[10pt,a4paper]{amsart}
\usepackage{amsmath,amssymb,amsthm,graphicx}
\usepackage{verbatim}
\newenvironment{reference}{\comment}{\endcomment}
%\newenvironment{reference}{}{}
\newenvironment{slogan}{\comment}{\endcomment}
\newenvironment{history}{\comment}{\endcomment}

% For commutative diagrams we use Xy-pic
\usepackage[all]{xy}

% We use 2cell for 2-commutative diagrams.
\xyoption{2cell}
% Theorem environments.
%
\theoremstyle{plain}
\newtheorem{theorem}[subsection]{Theorem}
\newtheorem{proposition}[subsection]{Proposition}
\newtheorem{lemma}[subsection]{Lemma}

\theoremstyle{definition}
\newtheorem{definition}[subsection]{Definition}
\newtheorem{example}[subsection]{Example}
\newtheorem{exercise}[subsection]{Exercise}
\newtheorem{situation}[subsection]{Situation}

\theoremstyle{remark}
\newtheorem{remark}[subsection]{Remark}
\newtheorem{remarks}[subsection]{Remarks}

\numberwithin{equation}{subsection}

% Macros
%
\def\lim{\mathop{\mathrm{lim}}\nolimits}
\def\colim{\mathop{\mathrm{colim}}\nolimits}
\def\Spec{\mathop{\mathrm{Spec}}}
\def\Hom{\mathop{\mathrm{Hom}}\nolimits}
\def\Ext{\mathop{\mathrm{Ext}}\nolimits}
\def\SheafHom{\mathop{\mathcal{H}\!\mathit{om}}\nolimits}
\def\SheafExt{\mathop{\mathcal{E}\!\mathit{xt}}\nolimits}
\def\Sch{\mathit{Sch}}
\def\Mor{\mathop{\mathrm{Mor}}\nolimits}
\def\Ob{\mathop{\mathrm{Ob}}\nolimits}
\def\Sh{\mathop{\mathit{Sh}}\nolimits}
\def\NL{\mathop{N\!L}\nolimits}
\def\CH{\mathop{\mathrm{CH}}\nolimits}
\def\proetale{{pro\text{-}\acute{e}tale}}
\def\etale{{\acute{e}tale}}
\def\QCoh{\mathit{QCoh}}
\def\Ker{\mathop{\mathrm{Ker}}}
\def\Im{\mathop{\mathrm{Im}}}
\def\Coker{\mathop{\mathrm{Coker}}}
\def\Coim{\mathop{\mathrm{Coim}}}

% Boxtimes
%
\DeclareMathSymbol{\boxtimes}{\mathbin}{AMSa}{"02}

%
% Macros for moduli stacks/spaces
%
\def\QCohstack{\mathcal{QC}\!\mathit{oh}}
\def\Cohstack{\mathcal{C}\!\mathit{oh}}
\def\Spacesstack{\mathcal{S}\!\mathit{paces}}
\def\Quotfunctor{\mathrm{Quot}}
\def\Hilbfunctor{\mathrm{Hilb}}
\def\Curvesstack{\mathcal{C}\!\mathit{urves}}
\def\Polarizedstack{\mathcal{P}\!\mathit{olarized}}
\def\Complexesstack{\mathcal{C}\!\mathit{omplexes}}
% \Pic is the operator that assigns to X its picard group, usage \Pic(X)
% \Picardstack_{X/B} denotes the Picard stack of X over B
% \Picardfunctor_{X/B} denotes the Picard functor of X over B
\def\Pic{\mathop{\mathrm{Pic}}\nolimits}
\def\Picardstack{\mathcal{P}\!\mathit{ic}}
\def\Picardfunctor{\mathrm{Pic}}
\def\Deformationcategory{\mathcal{D}\!\mathit{ef}}
\usepackage{cleveref}
\crefname{originalreference}{original reference}{original references}
\Crefname{originalreference}{Original reference}{Original references}
\makeatletter
\newlabel{definition-approximation-holds}{{08EM}{source}{}{sourcepointer-1}{}}
\newlabel{definition-approximation-holds@cref}{{[originalreference][0][]08EM}{[1][0][]source}{}{}{}}
\makeatother
\crefname{originalequation}{original equation}{original equations}
\Crefname{originalequation}{Original equation}{Original equations}
\makeatletter
\newlabel{equation-epsilon}{{08H7}{source}{}{sourcepointer-2}{}}
\newlabel{equation-epsilon@cref}{{[originalequation][0][]08H7}{[1][0][]source}{}{}{}}
\makeatother
\crefname{originalreference}{original reference}{original references}
\Crefname{originalreference}{Original reference}{Original references}
\makeatletter
\newlabel{lemma-derived-quasi-coherent-small-etale-site}{{071Q}{source}{}{sourcepointer-3}{}}
\newlabel{lemma-derived-quasi-coherent-small-etale-site@cref}{{[originalreference][0][]071Q}{[1][0][]source}{}{}{}}
\makeatother
\crefname{originalreference}{original reference}{original references}
\Crefname{originalreference}{Original reference}{Original references}
\makeatletter
\newlabel{lemma-descend-perfect}{{08HG}{source}{}{sourcepointer-4}{}}
\newlabel{lemma-descend-perfect@cref}{{[originalreference][0][]08HG}{[1][0][]source}{}{}{}}
\makeatother
\crefname{originalreference}{original reference}{original references}
\Crefname{originalreference}{Original reference}{Original references}
\makeatletter
\newlabel{lemma-descend-pseudo-coherent}{{08HE}{source}{}{sourcepointer-5}{}}
\newlabel{lemma-descend-pseudo-coherent@cref}{{[originalreference][0][]08HE}{[1][0][]source}{}{}{}}
\makeatother
\crefname{originalreference}{original reference}{original references}
\Crefname{originalreference}{Original reference}{Original references}
\makeatletter
\newlabel{lemma-epsilon-flat}{{08H8}{source}{}{sourcepointer-6}{}}
\newlabel{lemma-epsilon-flat@cref}{{[originalreference][0][]08H8}{[1][0][]source}{}{}{}}
\makeatother
\crefname{originallemma}{original lemma}{original lemmas}
\Crefname{originallemma}{Original lemma}{Original lemmas}
\makeatletter
\newlabel{perfect-lemma-approximation-affine}{{08EQ}{source}{}{sourcepointer-7}{}}
\newlabel{perfect-lemma-approximation-affine@cref}{{[originallemma][0][]08EQ}{[1][0][]source}{}{}{}}
\makeatother
\title{Approximation by perfect complexes over affine spaces}
\author{The Stacks Project Authors}
\date{stacks-project:spaces-perfect.tex:local-snapshot-20261012; GNU FDL 1.2 or later}
\begin{document}
\maketitle
\noindent\textit{Editorial note.} All delivered mathematical text is the author's own, copied from the pinned Stacks Project chapter file spaces-perfect.tex; no statement or proof step was rewritten, reordered or removed.\par\medskip
\noindent\textit{Editorial note.} The counted proof is the author's complete argument: the proof descends E along epsilon, transfers pseudo-coherence and support by the stalk comparison, applies the scheme approximation lemma, and pulls the perfect approximation back using descent of perfectness\par\medskip
\noindent\textit{Editorial note.} The label lemma-approximation-affine names the counted statement, and the proof environment immediately following it in the source is delivered verbatim as its proof.\par\medskip
\noindent\textit{Editorial note.} The counted statement is carried by the author's own lemma environment in the chapter source.\par\medskip
\noindent\textit{Allowed standard prerequisites (editorial).}\begin{itemize}
\item the definition of approximation for triples (T, E, m) (definition-approximation-holds of the same chapter)
\item the comparison of D\_QCoh on the small etale site with the Zariski site via epsilon (lemma-derived-quasi-coherent-small-etale-site of the same chapter)
\item descent of m-pseudo-coherence along epsilon (lemma-descend-pseudo-coherent of the same chapter) and the stalk comparison of cohomology sheaves (as in the proof of lemma-epsilon-flat)
\item the affine scheme approximation lemma (Derived Categories of Schemes, Lemma perfect-lemma-approximation-affine)
\item descent of perfectness along epsilon (lemma-descend-perfect of the same chapter)
\end{itemize}\medskip
\ifdefined\hypertarget\hypertarget{sourcepointer-1}{}\fi
\noindent\textit{Source reference.} Original reference 08EM: sources/stacks-spaces-perfect/pinned-original-source.tex, source lines 3575--3575. Prerequisite pointer invoked at proof line(s) 3654; it names a standard fact of the same chapter that the counted argument uses, and every step of the counted proof built on it is retained verbatim in the delivered body.\par\medskip
\ifdefined\hypertarget\hypertarget{sourcepointer-2}{}\fi
\noindent\textit{Source reference.} Original equation 08H7: sources/stacks-spaces-perfect/pinned-original-source.tex, source lines 320--320. Prerequisite pointer invoked at proof line(s) 3668; it names a standard fact of the same chapter that the counted argument uses, and every step of the counted proof built on it is retained verbatim in the delivered body.\par\medskip
\ifdefined\hypertarget\hypertarget{sourcepointer-3}{}\fi
\noindent\textit{Source reference.} Original reference 071Q: sources/stacks-spaces-perfect/pinned-original-source.tex, source lines 402--402. Prerequisite pointer invoked at proof line(s) 3670; it names a standard fact of the same chapter that the counted argument uses, and every step of the counted proof built on it is retained verbatim in the delivered body.\par\medskip
\ifdefined\hypertarget\hypertarget{sourcepointer-4}{}\fi
\noindent\textit{Source reference.} Original reference 08HG: sources/stacks-spaces-perfect/pinned-original-source.tex, source lines 3288--3288. Prerequisite pointer invoked at proof line(s) 3678; it names a standard fact of the same chapter that the counted argument uses, and every step of the counted proof built on it is retained verbatim in the delivered body.\par\medskip
\ifdefined\hypertarget\hypertarget{sourcepointer-5}{}\fi
\noindent\textit{Source reference.} Original reference 08HE: sources/stacks-spaces-perfect/pinned-original-source.tex, source lines 3133--3133. Prerequisite pointer invoked at proof line(s) 3672; it names a standard fact of the same chapter that the counted argument uses, and every step of the counted proof built on it is retained verbatim in the delivered body.\par\medskip
\ifdefined\hypertarget\hypertarget{sourcepointer-6}{}\fi
\noindent\textit{Source reference.} Original reference 08H8: sources/stacks-spaces-perfect/pinned-original-source.tex, source lines 330--330. Prerequisite pointer invoked at proof line(s) 3674; it names a standard fact of the same chapter that the counted argument uses, and every step of the counted proof built on it is retained verbatim in the delivered body.\par\medskip
\ifdefined\hypertarget\hypertarget{sourcepointer-7}{}\fi
\noindent\textit{Source reference.} Original lemma 08EQ: sources/stacks-chow/tags.txt, source lines 10681--10681. Prerequisite pointer invoked at proof line(s) 3676; it names a standard fact of a sibling chapter that the counted argument uses, and every step of the counted proof built on it is retained verbatim in the delivered body.\par\medskip
\begin{lemma}
\label{lemma-approximation-affine}
Let $S$ be a scheme. Let $X$ be an algebraic space over $S$ which is
representable by an affine scheme. Then approximation holds for every
triple $(T, E, m)$ as in Definition \ref{definition-approximation-holds}
such that there exists an integer $r \geq 0$ with
\begin{enumerate}
\item $E$ is $m$-pseudo-coherent,
\item $H^i(E)$ is supported on $T$ for $i \geq m - r + 1$,
\item $X \setminus T$ is the union of $r$ affine opens.
\end{enumerate}
In particular, approximation by perfect complexes holds for affine schemes.
\end{lemma}

\begin{proof}
Let $X_0$ be an affine scheme representing $X$. Let $T_0 \subset X_0$
by the closed subset corresponding to $T$. Let
$\epsilon : X_\etale \to X_{0, Zar}$ be the morphism
(\ref{equation-epsilon}). We may write $E = \epsilon^*E_0$ for some object
$E_0$ of $D_\QCoh(\mathcal{O}_{X_0})$, see
Lemma \ref{lemma-derived-quasi-coherent-small-etale-site}.
Then $E_0$ is $m$-pseudo-coherent, see
Lemma \ref{lemma-descend-pseudo-coherent}.
Comparing stalks of cohomology sheaves (see proof of
Lemma \ref{lemma-epsilon-flat})
we see that $H^i(E_0)$ is supported on $T_0$ for $i \geq m - r + 1$. By
Derived Categories of Schemes, Lemma \ref{perfect-lemma-approximation-affine}
there exists an approximation $P_0 \to E_0$ of
$(T_0, E_0, m)$. By Lemma \ref{lemma-descend-perfect}
we see that $P = \epsilon^*P_0$ is a perfect object of $D(\mathcal{O}_X)$.
Pulling back we obtain an approximation
$P = \epsilon^*P_0 \to \epsilon^*E_0 = E$ as desired.
\end{proof}

\end{document}
